\section{Individual alignment versus group alignment} \label{S2}

Distinguishing between the properties of individuals and organisations helps us understand the risks associated with ACI and reason about safety incidents like OAI-HF. In general, we rarely assume that organisations are the sum of their individual members. Instead, they have collective properties that stem from interactions among individuals, the social hierarchies that develop within them, and the social norms that the organisation develops over time \parencite{durkheim1895}. These collective properties can lead well-meaning people to staff institutions that cause harm. For example, before the 1986 Challenger disaster, no single NASA engineer sought to endanger an astronaut’s life \parencite{vaughan1996}. However, group pressures within NASA to normalise warning signs allowed departures from expected performance and, ultimately, culminated in the death of Challenger’s crew \parencite{vaughan1996}.

This insight underscores the risks of ACI and the limits of current model-alignment techniques. Fundamentally, we cannot assume that aligning an individual model ensures a safe collective when its agents collaborate. In April 2026, \textcite{shen_ai_2026} demonstrated this gap empirically. The authors copied a single aligned model to create an organisation and tasked the organisation with achieving 12 business-related tasks. They also gave a single copy of the model the same tasks to complete alone. Across all settings, \textcite{shen_ai_2026} judged that the organisations produced better business outcomes but worse ethical decisions than the individual. Their findings support an emerging consensus within the safety community: that model alignment does not generalise to multi-agent settings, and groups can become misaligned through interaction and collaboration \parencite{hammond2025}.

\citeauthor{shen_ai_2026}’s \parencite*{shen_ai_2026} observations expose the limitations of current safety techniques, such as supervised fine-tuning (SFT), reinforcement learning from human feedback (RLHF), and chain-of-thought (CoT) monitoring \parencite{baker2025cot,korbak2025cot, wei2022flan,ouyang2022,christiano2017}. These techniques operate at an individual level and therefore cannot determine whether a given response will combine with others to produce a harmful output, since each agent can meet the requirements set for it while the group’s combined output fails the requirements set for the whole \parencite{altmann_emergence_2024}. Imposing individual guardrails also fails to remove compounding biases within agentic collectives, where small systematic errors are repeated back to the models and so reinforced and amplified by the group \parencite{flint_group_2026, shumailov2024}. Finally, if a collective implicitly rewards or punishes agents based on effective performance, aligning agents to achieve individually specified objectives encourages them to dispense with guardrails to succeed within the collective \parencite{hammond_multi-agent_2025}. Thus, individual alignment is not compositional; to reach safe ACI, we must turn to group-level techniques.

\subsection{Defining group alignment}

What does it mean for a collective of agents to be aligned with human values and intentions? Is the group aligned only if every agent dogmatically follows rules for how to engage with each other and humans? Or can alignment be reflected only in the collective’s social norms and self-regulation mechanisms?

Here, it helps to consider distinct failures that researchers have associated with collective misalignment \parencite{hammond_multi-agent_2025,pierucci_institutional_2026}. Some researchers associate this term with a population settling on a position that its constituents would not reach individually, noting this as distinct from misalignment with human values \parencite{flint_group_2026}. Others measure collective misalignment against an external reference, such as a distribution of human safety standards or a published constitution \parencite{wang_devil_2026, shen_ai_2026}. Consequently, these separate but compatible positions imply that a definition of group alignment should support claims about whether a system is both safe for and supportive of human prosperity, while allowing benign drift from its agents’ priors. It must also be non-compositional and capture the difference between a collective of individually aligned agents and an aligned group. Informed by these ideas, we propose the following definition of group alignment:

\begin{quote}
\textit{A group of agents is group-aligned when its interaction structure, including the shared conventions that emerge between agents, reliably produces joint behaviour that respects human values.}
\end{quote}

Joint behaviour is an important term in this definition, since without it, every agent could satisfy its own specification while the group violates a global specification (\textit{composition}). Likewise, the group’s conventions must continue to respect human values under sustained pressure to abandon them, including pressure from within the group (\textit{durability}). Finally, across different pressures that the group might face, its interaction structure must not make violating human values the preferred option for an agent (\textit{incentive compatibility)}.

Under our definition, individual alignment is thus neither necessary nor sufficient for group alignment. A group of aligned agents can fail its three conditions, and equally, the group’s joint behaviour could be considered aligned even if it includes agents who, individually, are misaligned. Satisfying our definition requires only that a group settle into shared conventions that match human values, intentions, and ethical standards. The conventions hold because agents share them rather than individual agents prefer them, and they emerge through self-regulation rather than exogenously designed social structures.